\documentclass[10pt,a4paper]{amsart}
\usepackage[margin=19mm]{geometry}
\usepackage{amsmath,amssymb,mathtools}
\usepackage[scr=rsfs,cal=euler]{mathalfa}
\usepackage{xfrac}
\usepackage[unicode,hidelinks,bookmarks=false]{hyperref}
\newcommand{\ie}{i.e.\ }
\newcommand{\cf}{cf.\ }
\newcommand{\resp}{respectively\ }
\newcommand{\Subsection}{\subsection}
\providecommand{\nobreakdash}{\nobreak\hbox{-}}
\setlength{\emergencystretch}{2em}
\multlinegap0pt
\newcommand{\Res}{\mathrm{Res}}
\newcommand{\Disc}{\mathrm{Disc}}
\newcommand{\UU}{\mathbb{U}}
\newcommand{\ZZ}{\mathbb{Z}}
\newtheorem{theorem}{Theorem}
\newtheorem{lemma}{Lemma}
\title[Theorem 3.2 of arXiv:2508.21713]{Resultant of an equivariant polynomial system with respect to direct product of symmetric groups: the discriminant decomposition formula (Theorem 3.2)}
\author{Sonagnon Julien Owolabi, Ibrahim Nonkane, and Joel Tossa}
\date{Source: arXiv:2508.21713v1 (CC BY 4.0)}
\begin{document}
\maketitle
\noindent\emph{Source and licence.} Resultant of an equivariant polynomial system with respect to direct product of symmetric groups. Version arXiv:2508.21713v1 (CC BY 4.0), distributed under the Creative Commons Attribution 4.0 International licence, http://creativecommons.org/licenses/by/4.0/, as recorded in the arXiv OAI licence metadata for this exact version and shown on the abstract page https://arxiv.org/abs/2508.21713v1. Full rights notice: licenses/arxiv-2508.21713-CC-BY-4.0.txt.\par\smallskip
\noindent\textit{Editorial scope.} Counted unit: Theorem 3.2 of the paper (native source0.tex lines 743--788), the decomposition formula for the discriminant of a homogeneous polynomial invariant under the direct product of two symmetric groups, delivered together with the authors' own complete proof of it (source0.tex lines 790--792) as the single primary proof body. No mathematics is written, repaired or re-derived: statement and proof are the authors' own text line for line, in the authors' own notation ($\Res$, $\Disc$, $\mathcal{S}_{\{1,\ldots,p\}}$, $m_\lambda$, $\mu_1$, $\mu_2$, $r_1$, $r_2$); in particular the printed slips of the source are kept literal, including the full stop printed on its own line after ``Let's put $q=n-p$, $\Lambda=(\lambda,\lambda')\vdash(p,q)$'' in the retained Theorem 2.1, the colon that stands in place of a full stop at the end of every case heading in that theorem, and the two occurrences of an upper-case $F$ (namely $\Disc(F)$ and $F_{(\lambda,\lambda')}^{\{p+1\}}$) in the printed fourth case of the counted statement, where the surrounding text writes $f$; nothing was corrected. Required context retained uncounted, in the paper's own order: (a) lines 582--639, the paper's Theorem 2.1 (the general resultant decomposition formula) together with the authors' paragraph ``Idea of the proof'', which states that the detailed proof is deferred to an extended version of the paper -- this is retained precisely because the counted proof is a specialisation of that theorem; (b) lines 694--741, the opening of the paper's Section 3, the definition $d^{a(n,d)}\Disc(f)=\Res(f^{\{1\}},\ldots,f^{\{n\}})$ with $a(n,d)=((d-1)^n-(-1)^n)/d$, the authors' Lemma 3.1 with its complete proof (the partial derivatives of a $\mathcal{S}_{\{1,\ldots,p\}}\times\mathcal{S}_{\{p+1,\ldots,n\}}$-invariant homogeneous polynomial form an equivariant system) and the sentence that introduces the counted theorem. Omitted from this package and not redistributed: the paper's front matter and the rest of the source (all lines other than the four spans above, including the authors' Examples 2.2/3.3 and the paper's bibliography), and the paper's own preamble, which declares many macros the retained spans do not use; the wrapper re-states verbatim only the four notation macros the retained spans need ($\Res$, $\Disc$, $\UU$, $\ZZ$) and defines only the two theorem environments used ($theorem$, $lemma$). The wrapper numbers its own sections and theorem environments, so the retained Theorem 2.1 and Lemma 3.1 are renumbered; the paper's printed identifiers are named in this scope and in the item metadata. No retained span contains a graphics-inclusion command and the paper has no figures at all, so no image asset is shipped or declared.
\section*{Retained context: the paper's Theorem 2.1 (the decomposition formula) and the authors' ``Idea of the proof''}
\begin{theorem}\label{thm:maintheorem} Assume that $n\geq 2$ and  assume a system of $n$ homogeneous polynomials $f^{\{1\}},\ldots, f^{\{p\}}, f^{\{p+1\}},\ldots, f^{\{n\}}$ in $R[x_1,\ldots,x_n]$ of the same degree $d$ equivariant with respect to the direct product of two symmetric groups  $\mathcal{S}_{\{1
,\ldots,p\}}\times\mathcal{S}_{\{p+1
,\ldots,n\}} $ with $1\leq p<n$. Let's put $q=n-p$, $\Lambda=(\lambda, \lambda')\vdash (p, q)$. 
. 

\noindent $\bullet$ If $p\leq d$ and $q\leq d$ then :
\begin{multline*}
\Res\left(f^{\{1\}},\ldots,f^{\{p\}},f^{\{p+1\}},\ldots,f^{\{n\}}\right)=\\ \pm
 \prod_{\substack{\Lambda\vdash (p,q)}}
\Res\left( f_{\Lambda}^{\{1\}}, f_{\Lambda}^{\{1,2\}}, \ldots, f_{\Lambda}^{\{1,2,\ldots,r_1\}},f_{\Lambda}^{\{p+1\}}, f_{\Lambda}^{\{p+1,p+2\}}, \ldots, f_{\Lambda}^{\{p+1,p+2,\ldots,p+r_2\}} 
\right)^{m_{\lambda}m_{\lambda'}}.
\end{multline*}


\noindent $\bullet$ If $p> d$ and $q\leq d$ then
\begin{multline*}
\Res\left(f^{\{1\}},\ldots,f^{\{p\}},f^{\{p+1\}},\ldots,f^{\{n\}}\right)=
 \pm \left(f^{\{1,\ldots,d+1\}}\right)^{\mu_1}\\ \times\prod_{\substack{\Lambda\vdash (p,q)\\ r_1\leq d}}
\Res\left( f_{\Lambda}^{\{1\}}, f_{\Lambda}^{\{1,2\}}, \ldots, f_{\Lambda}^{\{1,2,\ldots,r_1\}},f_{\Lambda}^{\{p+1\}}, f_{\Lambda}^{\{p+1,p+2\}}, \ldots, f_{\Lambda}^{\{p+1,p+2,\ldots,p+r_2\}} 
\right)^{m_{\lambda}m_{\lambda'}}.
\end{multline*}


\noindent $\bullet$ If $p\leq d$ and $q> d$ then
\begin{multline*}
\Res\left(f^{\{1\}},\ldots,f^{\{p\}},f^{\{p+1\}},\ldots,f^{\{n\}}\right)=
 \pm \left(f^{\{p+1,\ldots,p+1+d\}}\right)^{\mu_2}\\ \times\prod_{\substack{\Lambda\vdash (p,q)\\ r_2\leq d}}
\Res\left( f_{\Lambda}^{\{1\}}, f_{\Lambda}^{\{1,2\}}, \ldots, f_{\Lambda}^{\{1,2,\ldots,r_1\}},f_{\Lambda}^{\{p+1\}}, f_{\Lambda}^{\{p+1,p+2\}}, \ldots, f_{\Lambda}^{\{p+1,p+2,\ldots,p+r_2\}} 
\right)^{m_{\lambda}m_{\lambda'}}.
\end{multline*}

\noindent $\bullet$ If $p> d$ and $q> d$ then
\begin{multline*}
\Res\left(f^{\{1\}},\ldots,f^{\{p\}},f^{\{p+1\}},\ldots,f^{\{n\}}\right)=
 \pm \left(f^{\{1,\ldots,d+1\}}\right)^{\mu_1} \times \left(f^{\{p+1,\ldots,p+1+d\}}\right)^{\mu_2}\\ \times\prod_{\substack{\Lambda\vdash (p,q)\\ r_1\leq d,r_2\leq d}}
\Res\left( f_{\Lambda}^{\{1\}}, f_{\Lambda}^{\{1,2\}}, \ldots, f_{\Lambda}^{\{1,2,\ldots,r_1\}},f_{\Lambda}^{\{p+1\}}, f_{\Lambda}^{\{p+1,p+2\}}, \ldots, f_{\Lambda}^{\{p+1,p+2,\ldots,p+r_2\}} 
\right)^{m_{\lambda}m_{\lambda'}}.
\end{multline*}



where
\begin{equation*}
   \mu_1:=nd^{n-1}-\sum_{ \substack{\lambda\vdash p \\ r_1 \leq d}} m_{\lambda} 
  \left(\sum_{j=1}^{r_1}\dfrac{d(d-1)\cdots(d-r_1+1)\times d^q}{d-j+1}+d(d-1)\cdots(d-r_1+1)\times qd^{q-1} \right).
  \end{equation*}
  and
\begin{equation*}
   \mu_2:=nd^{n-1}-\sum_{ \substack{\lambda'\vdash q \\ r_2 \leq d}} m_{\lambda'} 
  \left(\sum_{j=1}^{r_2}\dfrac{d(d-1)\cdots(d-r_2+1)\times d^p}{d-j+1}+d(d-1)\cdots(d-r_2+1)\times pd^{p-1} \right).
  \end{equation*}

\end{theorem}

\paragraph*{\it {\bf Idea of the proof}}
The main idea of the proof is the same as the one in \cite{BuKa16}. In fact It is clear that the system $\{f^{\{1\}},\ldots, f^{\{p\}} \}$ is equivariant respect to  $\mathcal{S}_{\{1,\ldots,p\}} $ and the system $\{ f^{\{p+1\}},\ldots f^{\{n\}} \}$ is equivariant with respect to $\mathcal{S}_{\{p+1,\ldots,n\}}$. The proof goes on by splitting the resultant of  $f^{{\{i\}}}$'s into several factors by means of their divided differences associated to  $\mathcal{S}_{\{1,\ldots,p\}} $ and respectively. For the rest of the proof , we mimick  the proof of \cite[Theorem 3.3]{BuKa16}. Indeed this is a generalization of the proof of \cite[Theorem 3.3]{BuKa16}. \\
%We begin by splitting the resultant of the  $F^{{\{i\}}}$'s into several factors by means of their divided differences (see paragraph 2.1.1 of \cite{BuKa16}). It is clear that the system $\{F^{\{1\}},\ldots, F^{\{p\}} \}$ is equivariant respect to  $\mathcal{S}_{\{1,\ldots,p\}} $ and the system $\{ F^{\{p+1\}},\ldots F^{\{n\}} \}$ is equivariant with respect to $\mathcal{S}_{\{p+1,\ldots,n\}}$.
For the sake of the number of pages,the detailed discussions of the proof  will be published later in an extended version of this paper. \\

\section*{Retained context: the discriminant identity and Lemma 3.1 of the paper (paper \S3)}
\section{Discriminant of a homogeneous polynomial invariant under direct product of symmetric groups }\label{sec:discriminant}

In this section, we will use the previous theorem  Theorem {\ref{thm:maintheorem}} to  develop  a decomposition formula for the discriminant of an invariant homogeneous polynomial under the action of  $\mathcal{S}_{\{1
,\ldots,p\}}\times\mathcal{S}_{\{p+1
,\ldots,n\}}.$ 
%Let $R$ be commutative ring and denote by $R[x_1,\ldots,x_n]$ the ring of polynomials in $n \geq 2$ variables which is graded with the usual weights: $\deg (x_i)=1$ for all $i \in \{ 1,\ldots,n\}.$. 
Let $f \in R[x_1,\ldots,x_p,x_{p+1},\ldots,x_n]$ of degree $d$ be a homogeneous polynomial that is invariant under direct product  $\mathcal{S}_{\{1
,\ldots,p\}}\times\mathcal{S}_{\{p+1
,\ldots,n\}} $ of symmetric groups with $1\leq p< n$.

For all $\sigma_1\in \mathcal{S}_{\{1
,\ldots,p\}}$ and for all $\sigma_2\in \mathcal{S}_{\{p+1
,\ldots,n\}}$, we have :


\begin{eqnarray}
%\begin{equation}\label{eq:F1}
(\sigma_1,\sigma_2)\Big(f \Big)(x_1,\ldots,x_p,x_{p+1},\ldots,x_n) &=&f(x_{\sigma_1(1)},\ldots x_{\sigma_1(p)},x_{\sigma_2(p+1)},\ldots,x_{\sigma_2(n)}). \\
%\end{equation}
%\begin{equation}\label{eq:F}
%(\sigma_1,\sigma_2)\Big(F(x_1,\ldots,x_p,x_{p+1},\ldots,x_n)\Big)
 &=& f(x_1,\ldots,x_p,x_{p+1},\ldots,x_n)
%\end{equation}
\end{eqnarray}

We  will denote the partial derivatives of $F$ by
$$f^{\{i\}}(x_1,\ldots,x_p,x_{p+1},\ldots,x_n):=\frac{\partial f}{\partial x_{i}} (x_1,\ldots,x_p,x_{p+1},\ldots,x_n), \ i=1,\ldots n .$$ 
The discriminant of $F$ is defined by the equality
\begin{equation}\label{eq:discressym}
d^{a(n,d)}\Disc(f)=\Res\left(  f^{\{1\}},f^{\{2\}},\ldots,f^{\{n\}} \right) \in \UU
\end{equation} where 
$$a(n,d):=\frac{(d-1)^{n}-(-1)^{n}}{d} \in \ZZ.$$
and that it is homogeneous of degree $n(d-1)^{n-1}$.

\begin{lemma}
The set $\{ f^{\{1\}},f^{\{2\}},\ldots,f^{\{p\}},f^{\{p+1\}},\ldots,f^{\{n\}} \}$ of partial derivatives of a  $\mathcal{S}_{\{1,\ldots,p\}}\times\mathcal{S}_{\{p+1,\ldots,n\}} $-invariant homogeneous polynomial $f$ is is an equivariant polynomial system with respect to $\mathcal{S}_{\{1,\ldots,p\}}\times\mathcal{S}_{\{p+1,\ldots,n\}} $.
\end{lemma}
\begin{proof}
We will use the canonical inclusions  $ \mathcal{S}_{\{1,\ldots,p\}} \to \mathcal{S}_{\{1,\ldots,p\}}\times\mathcal{S}_{\{p+1,\ldots,n\}}, \sigma_1 \mapsto (\sigma_1, e_2)$ and  $ \mathcal{S}_{\{p+1,\ldots,n\}} \to \mathcal{S}_{\{1,\ldots,p\}}\times\mathcal{S}_{\{p+1,\ldots,n\}}, \sigma_2 \mapsto (e_1,\sigma_2)$ where $e_1, e_2$ are unit elements of $ \mathcal{S}_{\{1,\ldots,p\}}$ and $\mathcal{S}_{\{p+1,\ldots,n\}}$ respectively.
For all $i\in \{1,\ldots,p\}$ and $\sigma_1 \in \mathcal{S}_{\{1,\ldots,p\}},$ we have $\sigma_1\Big(f^{\{i\}}\Big)=\sigma_1\Big(\frac{\partial f}{\partial x_{i}}\Big)=\frac{\partial (\sigma_1f)}{\partial x_{\sigma_1(i)}}=\frac{\partial f}{\partial x_{\sigma_1(i)}} =f^{\{ {\sigma_1(i)} \}}.$ For all $j\in \{p+1,\ldots,n\}$ and $\sigma_2 \in \mathcal{S}_{\{p+1,\ldots,n\}},$ , $\sigma_2\Big(f^{\{j\}}\Big)=\sigma_2\Big(\frac{\partial f}{\partial x_{j}}\Big)=\frac{\partial (\sigma_2 f)}{\partial x_{\sigma_2(j)}}=\frac{\partial f}{\partial x_{\sigma_2(j)}} =f^{\{ {\sigma_2(j)} \}}.$ For all $k \in \{1,\ldots,p\}$ $\sigma= (\sigma_1, \sigma_2) \in \mathcal{S}_{\{1,\ldots,p\}}\times\mathcal{S}_{\{p+1,\ldots,n\}}$, we have $\sigma \Big( f^{\{k\}} \Big)= \begin{cases}
\sigma_1 \Big( f^{\{k\}} \Big)= f^{\{ {\sigma_1(k)} \}}\  \mbox{if}\  k\in \{1,\ldots,p\} \\
\sigma_2 \Big( f^{\{k\}} \Big)= f^{\{ {\sigma_2(k)} \}}\  \mbox{if}\  k\in \{p+1,\ldots,n\}.
\end{cases}$
Hence $\sigma \Big( f^{\{k\}} \Big) \in  \{ f^{\{1\}},\ldots ,f^{\{n\}} \}$, for all  $k \in \{1,\ldots,p\}$ $\sigma= (\sigma_1, \sigma_2) \in \mathcal{S}_{\{1,\ldots,p\}}\times\mathcal{S}_{\{p+1,\ldots,n\}}$. and the set of partial derivative of $f$ form  an equivariant polynomial system with respect to $\mathcal{S}_{\{1,\ldots,p\}}\times\mathcal{S}_{\{p+1,\ldots,n\}} $.
\end{proof}


As a consequence of this lemma, Theorem \ref{thm:maintheorem} can be applied in order to decompose the resultant of the polynomials $f^{\{1\}},f^{\{2\}},\ldots,f^{\{p\}},f^{\{p+1\}},\ldots f^{\{n\}}$ and hence, by \eqref{eq:discressym}, to decompose the discriminant of the $\mathcal{S}_{\{1,\ldots,p\}}\times\mathcal{S}_{\{p+1,\ldots,n\}} $-invariant polynomial $f$. %We take again the notation of \S \ref{subsec:divdiff} and  \S \ref{subsec:mainthm}.

\section*{Counted claim: Theorem 3.2 of the paper, the discriminant decomposition formula, and the authors' complete proof of it}
\begin{theorem}\label{thm:discmaintheorem} Assume that $n\geq 2$ and $d\geq 2$. With the above notation, the following equalities hold.

\noindent $\bullet$ If $p< d$ and $q< d$ then :
\begin{multline*}
d^{a(n,d)}\Disc\left(f\right)=\\
 \prod_{\substack{\Lambda\vdash (p,q)}}
\Res\left( f_{(\lambda,\lambda')}^{\{1\}}, f_{(\lambda,\lambda')}^{\{1,2\}}, \ldots, f_{(\lambda,\lambda')}^{\{1,2,\ldots,r_1\}},f_{(\lambda,\lambda')}^{\{p+1\}}, f_{(\lambda,\lambda')}^{\{p+1,p+2\}}, \ldots, f_{(\lambda,\lambda')}^{\{p+1,p+2,\ldots,p+r_2\}} 
\right)^{m_{\lambda}m_{\lambda'}}.
\end{multline*}


\noindent $\bullet$ If $p\geqslant d$ and $q< d$ then :
\begin{multline*}
d^{a(n,d)}\Disc\left(f\right)=
  \left(f^{\{1,\ldots,d\}}\right)^{\mu_1}\\ \times\prod_{\substack{\Lambda\vdash (p,q)\\ r_1< d}}
\Res\left( f_{(\lambda,\lambda')}^{\{1\}}, f_{(\lambda,\lambda')}^{\{1,2\}}, \ldots, f_{(\lambda,\lambda')}^{\{1,2,\ldots,r_1\}},f_{(\lambda,\lambda')}^{\{p+1\}}, f_{(\lambda,\lambda')}^{\{p+1,p+2\}}, \ldots, f_{(\lambda,\lambda')}^{\{p+1,p+2,\ldots,p+r_2\}} 
\right)^{m_{\lambda}m_{\lambda'}}.
\end{multline*}

\noindent $\bullet$ If $p< d$ and $q\geqslant d$ then :
\begin{multline*}
d^{a(n,d)}\Disc\left(f\right)=
  \left(f^{\{p+1,\ldots,p+d\}}\right)^{\mu_2}\\ \times\prod_{\substack{\Lambda\vdash (p,q)\\ \ r_2< d}}
\Res\left( f_{(\lambda,\lambda')}^{\{1\}}, f_{(\lambda,\lambda')}^{\{1,2\}}, \ldots, f_{(\lambda,\lambda')}^{\{1,2,\ldots,r_1\}},f_{(\lambda,\lambda')}^{\{p+1\}}, f_{(\lambda,\lambda')}^{\{p+1,p+2\}}, \ldots, f_{(\lambda,\lambda')}^{\{p+1,p+2,\ldots,p+r_2\}} 
\right)^{m_{\lambda}m_{\lambda'}}.
\end{multline*}

\noindent $\bullet$ If $p\geqslant d$ and $q\geqslant d$ then :
\begin{multline*}
d^{a(n,d)}\Disc\left(F\right)=
  \left(f^{\{1,\ldots,d\}}\right)^{\mu_1}\times \left(f^{\{p+1,\ldots,p+d\}}\right)^{\mu_2}\\ \times\prod_{\substack{\Lambda\vdash (p,q)\\ r_1< d,\ r_2< d}}
\Res\left( f_{(\lambda,\lambda')}^{\{1\}}, f_{(\lambda,\lambda')}^{\{1,2\}}, \ldots, f_{(\lambda,\lambda')}^{\{1,2,\ldots,r_1\}},F_{(\lambda,\lambda')}^{\{p+1\}}, f_{(\lambda,\lambda')}^{\{p+1,p+2\}}, \ldots, f_{(\lambda,\lambda')}^{\{p+1,p+2,\ldots,p+r_2\}} 
\right)^{m_{\lambda}m_{\lambda'}}.
\end{multline*}

where
 \begin{equation*}
    \mu_1:=n(d-1)^{n-1}-\sum_{ \substack{\lambda\vdash p \\ r_1 < d}} m_{\lambda} 
   \left(\sum_{j=1}^{r_1}\dfrac{(d-1)\cdots(d-r_1)\times (d-1)^q}{d-j}+(d-1)\cdots(d-r_1)\times q(d-1)^{q-1} \right).
   \end{equation*}
   and
 \begin{equation*}
    \mu_2:=n(d-1)^{n-1}-\sum_{ \substack{\lambda'\vdash q \\ r_2 < d}} m_{\lambda'} 
   \left(\sum_{j=1}^{r_2}\dfrac{(d-1)\cdots(d-r_2)\times (d-1)^p}{d-j}+(d-1)\cdots(d-r_2)\times p(d-1)^{p-1} \right).
   \end{equation*}
\end{theorem}

\begin{proof}
These formulas are obtained by specialization of the formulas given in Theorem \ref{thm:maintheorem} with the difference that the polynomials $f^{\{i\}}$, $i=1,\ldots,n$ are of degree $d-1$ in our setting (and not of degree $d$ as in Theorem \ref{thm:maintheorem}).
\end{proof}

\begin{thebibliography}{99}
\bibitem[1]{BuKa16} Laurent Bus{\'e} and Anna Karasoulou. Resultant of an equivariant polynomial system with respect to the symmetric group. arXiv:1407.2799 v1[ math.AC] 10 Juillet 2014.
\end{thebibliography}
\end{document}
